\section*{Chapter \ref{ch:b3:palaeomagnetism} --- Geomagnetism and Palaeomagnetism}

\begin{solution}{exo:b3:palaeomagnetism:1}
$\tan I = 2\tan\ang{30} = 1.1547$, so $I = \ang{49.1}$. Conversely, $\tan\lambda
= \tfrac12\tan(\ang{-20}) = -0.1820$, so $\lambda = \ang{-10.3}$, that is
\ang{10.3}\,S.
\end{solution}

\begin{solution}{exo:b3:palaeomagnetism:2}
$H = \sqrt{20000^2 + 2000^2} = \qty{20100}{nT}$;
$F = \sqrt{20000^2 + 2000^2 + 40000^2} = \qty{44770}{nT}$. The magnetic
declination is $\arctan(-2000/20000) = \ang{-5.7}$, \ang{5.7} west of north;
the inclination $\arctan(40000/20100) = \ang{63.3}$, downwards.
\end{solution}

\begin{solution}{exo:b3:palaeomagnetism:3}
$B_0 = -g_1^0 = \qty{29350}{nT}$ at the equator, and $2B_0 = \qty{58700}{nT}$ at
the poles. $m = 4\pi a^3B_0/\mu_0 = a^3B_0/\num{e-7} = \num{2.5862e20}\times
\num{2.935e-5}/\num{e-7} = \qty{7.59e22}{A.m^2}$, a little less than the
\qty{7.69e22}{A.m^2} of the full dipole, which adds $g_1^1$ and $h_1^1$.
\end{solution}

\begin{solution}{exo:b3:palaeomagnetism:4}
(a) A \omterm{def:b3:palaeomagnetism:nrm}{thermoremanent magnetisation}, recording the field as the lava cooled
through the \omterm{def:b3:palaeomagnetism:blocking}{blocking temperatures} of its grains. (b) A \omterm{def:b3:palaeomagnetism:nrm}{detrital remanent
magnetisation}, recording the field within the time the silt took to settle and
be locked by compaction. (c) A \omterm{def:b3:palaeomagnetism:nrm}{chemical remanent magnetisation}, recording the
field during the growth of the haematite, during early burial and so somewhat
after deposition. (d) A viscous remanence along the present field, added to the
original TRM by the grains of shortest relaxation time; it is removed by the
first steps of demagnetisation.
\end{solution}

\begin{solution}{exo:b3:palaeomagnetism:5}
The unit vectors $(\cos I\cos D,\ \cos I\sin D,\ \sin I)$ are
\begin{gather*}
  (0.6545, 0.1391, 0.7431),\quad (0.6670, 0.2695, 0.6947),\\
  (0.5235, 0.1501, 0.8387),\quad (0.5985, 0.0629, 0.7986).
\end{gather*}
The sum is $(2.4435, 0.6216, 3.0751)$, of length $R = 3.9766$. Mean:
$D = \arctan(0.6216/2.4435) = \ang{14.3}$, $I = \arcsin(3.0751/3.9766) =
\ang{50.7}$. $k = 3/(4 - 3.9766) = 128$. With $20^{1/3} - 1 = 1.7144$,
$\cos\alpha_{95} = 1 - (0.0234/3.9766)\times1.7144 = 0.98991$, so
$\alpha_{95} = \ang{8.1}$: with four specimens the cone stays wide, however
tight the group.
\end{solution}

\begin{solution}{exo:b3:palaeomagnetism:6}
$\lambda = \arctan(\tfrac12\tan\ang{50}) = \ang{30.79}$, so $p = \ang{59.21}$.
$\sin\lambda_p = \sin\ang{40}\cos\ang{59.21} + \cos\ang{40}\sin\ang{59.21}
\cos\ang{20} = 0.9474$, $\lambda_p = \ang{71.34}$\,N. $\sin\beta =
\sin\ang{59.21}\sin\ang{20}/\cos\ang{71.34} = 0.9184$, $\beta = \ang{66.66}$.
Since $\cos p = 0.5119 < \sin\ang{40}\times0.9474 = 0.6090$, the second case
applies: $\varphi_p = 10 + 180 - 66.66 = \ang{123.34}$\,E. The pole is at
\ang{71.3}\,N, \ang{123.3}\,E.
\end{solution}

\begin{solution}{exo:b3:palaeomagnetism:7}
(a) $(7.69 - 8.32)/8.32 = \qty{-7.6}{\percent}$ in 125 years, or
\qty{-6.1}{\percent} per century. (b) $29350/12.6 = 2329$, about 2300 years.
It is no forecast: the rate itself changes from one five-year model to the
next, the moment has been both higher and lower in the past without
reversing, and the field reverses at irregular intervals that no rate
predicts.
\end{solution}

\begin{solution}{exo:b3:palaeomagnetism:8}
Between 20 and \qty{300}{\celsius} the vector lost $(6.4, 0.0, 13.6)$: a
component of \qty{15.0}{mA/m}, with $D = \ang{0.0}$ and
$I = \arctan(13.6/6.4) = \ang{64.8}$. Above \qty{300}{\celsius} it lost $(25.8, 18.1, 24.6)$:
\qty{40.0}{mA/m}, $D = \arctan(18.1/25.8) = \ang{35.1}$, $I =
\arctan(24.6/31.5) = \ang{38.0}$. The second, which goes to the origin, is the
ChRM; the first, along a northward, steep direction, removed at low
temperature, is a viscous overprint along the recent field.
\end{solution}

\begin{solution}{exo:b3:palaeomagnetism:9}
$\lambda = \arctan(\tfrac12\tan(\ang{-35})) = \ang{-19.3}$. For $I = \ang{-40}$
and \ang{-30}: \ang{-22.8} and \ang{-16.1}, an interval of \ang{6.7}.
$p = \ang{90} - \lambda = \ang{109.3}$, $\cos^2p = 0.1096$:
$\mathrm dp = 2.5\times(1 + 0.329) = \ang{3.3}$ and $\mathrm dm =
5\times\sin\ang{109.3}/\cos\ang{35} = \ang{5.8}$.
\end{solution}

\begin{solution}{exo:b3:palaeomagnetism:10}
(a) $\tan\lambda = \tfrac12$, $\lambda = \ang{26.6}$. $F$ is $B_0$ at the
equator and $2B_0$ at the poles: a ratio of 2. (b) $\tan I = 2\tan\lambda$
gives $(1 + \tan^2I)\,\mathrm dI = 2(1 + \tan^2\lambda)\,\mathrm d\lambda$, so
\[
  \frac{\mathrm dI}{\mathrm d\lambda} = \frac{2/\cos^2\lambda}{1 + 4\tan^2\lambda}
  = \frac{2}{\cos^2\lambda + 4\sin^2\lambda} = \frac{2}{1 + 3\sin^2\lambda}.
\]
An error $\delta I$ gives $\delta\lambda = \tfrac12(1 + 3\sin^2\lambda)\,\delta I$,
from $\tfrac12\delta I$ at the equator to $2\,\delta I$ at the poles: the
largest error is near the poles. (c) Near the equator the inclination changes
twice as fast as the latitude, so a given error of inclination costs half as
much latitude; near the poles the inclination hardly changes with latitude.
\end{solution}

\begin{solution}{exo:b3:palaeomagnetism:11}
(a) $\cos p = \sin(\ang{-35})\sin\ang{70} + \cos\ang{35}\cos\ang{70}\cos\ang{150}
= -0.5390 - 0.2426 = -0.7816$, $p = \ang{141.41}$. Then $\tan I_x = 2\cot p =
-2.506$, $I_x = \ang{-68.2}$, and $\lambda_x = \ang{90} - p = \ang{-51.4}$.
$\cos D_x = (\sin\ang{70} - \sin(\ang{-35})\cos p)/(\cos\ang{35}\sin p) =
0.9617$, and since the pole lies east of the site's meridian ($\sin\ang{150} >
0$), $D_x = \ang{15.9}$. (b) $\lambda_o = \arctan(\tfrac12\tan(\ang{-60})) =
\ang{-40.9}$. $R = 30.0 - 15.9 = \ang{14.1}$: a clockwise rotation of the block
about a vertical axis. $P = \lambda_x - \lambda_o = -51.4 + 40.9 =
\ang{-10.5}$: the formation lay \ang{10.5} further north than the stable plate
says it should, so the block has since moved about \ang{10.5}, some
\qty{1170}{km}, southwards with respect to the plate.
\end{solution}

\begin{solution}{exo:b3:palaeomagnetism:12}
(a) With $\mathrm d\Omega = 2\pi\sin\theta'\,\mathrm d\theta'$,
\[
  \int_0^\theta\frac{\kappa e^{\kappa\cos\theta'}}{4\pi\sinh\kappa}\,2\pi\sin\theta'\,
  \mathrm d\theta' = \frac{e^\kappa - e^{\kappa\cos\theta}}{2\sinh\kappa}
  = \frac{e^\kappa - e^{\kappa\cos\theta}}{e^\kappa - e^{-\kappa}} .
\]
(b) For large $\kappa$, $e^{-\kappa}$ is negligible; dividing by $e^\kappa$
leaves $1 - e^{-\kappa(1 - \cos\theta)} \approx 1 - e^{-\kappa\theta^2/2}$. It
equals $1 - e^{-1} = \qty{63}{\percent}$ when $\kappa\theta^2/2 = 1$, $\theta =
\sqrt{2/\kappa}$ radians, or $\ang{81}/\sqrt\kappa$. (c) $81/\sqrt{50} =
\ang{11.5}$; the exact formula gives \qty{63.1}{\percent}.
\end{solution}

\begin{solution}{pb:b3:palaeomagnetism:1}
\textbf{1.} Haematite in a red bed may grow from iron oxides and hydroxides
long after the sediment was laid down, giving a chemical remanence of that
later age. A positive fold test, if the beds were folded soon after
deposition, or a positive reversal test across closely spaced beds of both
polarities would show the remanence to be early.

\textbf{2.} $(\cos\ang{37}\cos\ang{50}, \cos\ang{37}\sin\ang{50}, \sin\ang{37})
= (0.5134, 0.6118, 0.6018)$.

\textbf{3.} $R = \sqrt{3.8943^2 + 2.6500^2 + 3.5717^2} = 5.9114$. $D =
\arctan(2.6500/3.8943) = \ang{34.2}$, $I = \arcsin(3.5717/5.9114) = \ang{37.2}$.

\textbf{4.} $k = 5/(6 - 5.9114) = 56.5$. With $20^{1/5} - 1 = 0.8206$,
$\cos\alpha_{95} = 1 - (0.0886/5.9114)\times0.8206 = 0.98771$, $\alpha_{95} =
\ang{9.0}$.

\textbf{5.} B: $k = 6/(7 - 6.8979) = 58.8$; $20^{1/6} - 1 = 0.6475$,
$\cos\alpha_{95} = 1 - (0.1021/6.8979)\times0.6475 = 0.99042$, $\alpha_{95} =
\ang{7.9}$. C: $k = 7/(8 - 7.8707) = 54.1$; $20^{1/7} - 1 = 0.5340$,
$\cos\alpha_{95} = 1 - (0.1293/7.8707)\times0.5340 = 0.99123$, $\alpha_{95} =
\ang{7.6}$.

\textbf{6.} B has the largest $k$, the tightest group; C has the smallest
$\alpha_{95}$, the best-defined mean, because it has the most sites. The three
are close: all three formations are well grouped.

\textbf{7.} $\lambda_A = \arctan(\tfrac12\tan\ang{37.17}) = \ang{20.76}$\,N,
$\lambda_B = \ang{32.15}$\,N, $\lambda_C = \ang{45.33}$\,N.

\textbf{8.} $p = 90 - 20.76 = \ang{69.24}$. $\sin\lambda_p =
\sin\ang{47}\cos\ang{69.24} + \cos\ang{47}\sin\ang{69.24}\cos\ang{34.23} =
0.2593 + 0.5272 = 0.7865$: $\lambda_p = \ang{51.86}$\,N.

\textbf{9.} $\sin\beta = \sin\ang{69.24}\sin\ang{34.23}/\cos\ang{51.86} =
0.8516$, $\beta = \ang{58.40}$. Since $\cos p = 0.3545 < \sin\ang{47}\times
0.7865 = 0.5752$, the second case applies, $\cos\beta < 0$:
$\varphi_p = 15.00 + 180 - 58.40 = \ang{136.60}$\,E.

\textbf{10.} $\cos^2p = 0.1257$: $\mathrm dp = 4.50\times1.377 = \ang{6.2}$;
$\mathrm dm = 8.99\times\sin\ang{69.24}/\cos\ang{37.17} = \ang{10.5}$.

\textbf{11.} From \ang{51.9}\,N at \qty{60}{Ma} to \ang{66.5}\,N at
\qty{30}{Ma} and \ang{87.1}\,N at \qty{4}{Ma}, on nearly the same meridian
(\ang{136}--\ang{139}\,E): a path towards the geographic pole, reached within
its cone by the youngest formation.

\textbf{12.} The ash bed, at \qty{22}{m}, lies in the second normal zone,
19.7--\qty{24.6}{m}; with an age near \qty{3.2}{Ma} that zone is C2An.2n
(3.207--3.116). Then the zone 4.2--\qty{14.8}{m} is C2An.3n, the zone
29.2--\qty{54.0}{m} C2An.1n, the reversed base belongs to the chron below
C2An, and the reversed top to the chron above it.

\textbf{13.} $(19.7 - 4.2)/(3.596 - 3.207) = 15.5/0.389 = \qty{39.8}{m/Myr}$;
$(54.0 - 24.6)/(3.116 - 2.581) = 29.4/0.535 = \qty{55.0}{m/Myr}$.

\textbf{14.} Base: $3.596 + 4.2/39.8 = \qty{3.70}{Ma}$. Top: $2.581 -
10.0/55.0 = \qty{2.40}{Ma}$.

\textbf{15.} Only reversal boundaries are dated, a few hundred thousand years
apart; between them the rate is assumed constant. A change of subsidence, of
sediment supply or of climate, or a gap in deposition, would change it, and a
hiatus could even remove a short zone and mislead the match.

\textbf{16.} $\cos p = \sin\ang{47}\sin\ang{78} + \cos\ang{47}\cos\ang{78}
\cos\ang{175} = 0.7154 - 0.1413 = 0.5741$, $p = \ang{54.96}$. $\tan I_x =
2\cot p = 1.4023$, $I_x = \ang{54.5}$; $\lambda_x = \ang{35.04}$\,N.

\textbf{17.} $\cos D_x = (\sin\ang{78} - \sin\ang{47}\times0.5741)/
(\cos\ang{47}\sin\ang{54.96}) = 0.99976$; the pole lies east of the site's
meridian ($\sin\ang{175} > 0$), so $D_x = \ang{1.3}$.

\textbf{18.} Expected \omterm{def:b3:palaeomagnetism:vgp}{palaeolatitudes} from $I_x$ = \ang{54.5}, \ang{60.0} and
\ang{63.6}: \ang{35.03}, \ang{40.89} and \ang{45.21}. $P_A = 35.03 - 20.76 =
\ang{14.3}$, $P_B = 40.89 - 32.15 = \ang{8.7}$, $P_C = 45.21 - 45.33 =
\ang{-0.1}$.

\textbf{19.} $R_A = 34.2 - 1.3 = \ang{32.9}$, $R_B = 23.9 - 2.9 = \ang{21.0}$,
$R_C = 3.4 - 2.5 = \ang{0.9}$, all clockwise.

\textbf{20.} $\Delta D_o = \arcsin(\sin\ang{8.99}/\cos\ang{37.17}) = \ang{11.3}$;
$\Delta D_x = \arcsin(\sin\ang{3.0}/\cos\ang{54.5}) = \ang{5.2}$; $\Delta R_A =
\sqrt{11.3^2 + 5.2^2} = \ang{12.4}$.

\textbf{21.} $R_A = \ang{32.9} \pm \ang{12.4}$ and $R_B = \ang{21.0} \pm
\ang{13.8}$ are significant; $R_C = \ang{0.9} \pm \ang{17.9}$ is not.

\textbf{22.} Sixty million years ago Kesta lay about \ang{14} of latitude
south of where the stable continent puts it, and was turned about \ang{33}
anticlockwise from its present orientation. By \qty{30}{Ma} it had come about
\ang{6} of those \ang{14} north and turned \ang{12} of those \ang{33}
clockwise; by \qty{4}{Ma} it had reached its present place on the continent,
with no rotation left. It drifted north along the continent's margin, turning
clockwise as it went, and was in place by \qty{4}{Ma}.

\textbf{23.} $47.00 - 20.76 = \ang{26.24}$, or $26.24\times111.19 =
\qty{2918}{km}$. This includes the motion of the stable continent itself;
relative to it, Kesta moved $P_A = \ang{14.3}$.

\textbf{24.} $I = \ang{28.18}$ and \ang{46.16} give $\lambda = \ang{15.00}$ and
\ang{27.50}: a half-width of \ang{6.25}, and one standard deviation of
$6.25/1.96 = \ang{3.19}$, or \qty{355}{km}. Since formation A, Kesta has drifted
northwards by
\textbf{\qty[text-series-to-math=true]{2918 \pm 355}{km}}.
\end{solution}
